\section{The certificate language}
\label{sec:certificates}

A claim about interaction is admitted into the catalog only as a certificate: a record
with one field for each thing the claim asserts, usually with an audit. The catalog has
twenty-three kinds of certificate, in four clusters. In every cluster validity has the
same shape: the fields the claim names are present, the Boolean flags it needs are set,
the stated implications between flags hold, and, where the record has one, the audit has
an entry. The supplement gives the validity predicate and the Lean fields of
every certificate (\Cref{supp:fields}).

\paragraph{Access and admission.} An access record gives a content seven Boolean
coordinates (\Cref{def:access}). An admission step is a typed transition of one of five
kinds; an admission regime records the routes by which a first extension may be
authorized. Neutral provisioning and prospective commitment records certify a source from
outside the system and a registration made before its evidence. A bridge contract records
what a transport between theories keeps, adds and loses; transport is licensed only by an
accepted, well-formed bridge, never by a citation.

\paragraph{Contact and join.} A contact surface and a contact witness record that two
theories touched (\Cref{def:contact}). A join entry records a candidate joint theory with
its parents, ledgers and one of seven statuses (\Cref{def:join-status}). The strict-join
certificate records the claim that a join is more than its parents
(\Cref{def:strict-certificate}). Three further records carry provenance and price. The
source-independence gate compares the ancestries of the two sides: each source traces
back to a root source, and the sides have \emph{disjoint roots} when their roots differ.
The cost ledger records what the join paid. An observer-occupancy record states how much
of an observer's or instrument's bounded capacity the interaction uses, its
\emph{occupancy}, and how much of that is charged to the ledger.

\paragraph{Enablement and descent.} An enablement record states that one theory makes an
operation available to another (\Cref{def:enablement}). A transmission record carries a
payload between layers. An operational profile and a negative-evidence record fix what a
rule's reachability and a failed search establish, and a coverage certificate is what a
claim of non-interaction must carry.

\paragraph{Dynamics and residue.} A route comparison records two orders of the same
members (\Cref{def:route}), a drive certificate is what an arrow claim needs
(\Cref{def:drive}), a residual ledger records what an interaction leaves undischarged, and a
contact-quantity record states what a claimed conservation includes.

The simplest record shows what a certificate looks like.

\begin{definition}[Access record]
\label{def:access}
An \emph{access record} for a content assigns seven Boolean coordinates: whether it is
\expressible, \present, \exposed, \recoverable, \admissible, \reachable{} and
\occurrent. The record is \emph{coherent} when
\[
  \occurrent\Rightarrow\reachable\Rightarrow\admissible\Rightarrow\present\Rightarrow
  \expressible,
  \qquad
  \recoverable\Rightarrow\exposed\Rightarrow\present,
\]
and its audit is non-empty.
\end{definition}

The implications run one way only, and the reference world contains a coherent record for
each failed converse (\Cref{sec:separations}). Of the $2^7=128$ assignments of the seven
coordinates, exactly 14 satisfy these two chains of implications, the
\emph{implication spine} of the record. A single record already distinguishes readings that ordinary usage runs
together. Consider a content that is \expressible, \present{} and \admissible{} but not
\exposed, \recoverable, \reachable{} or \occurrent. The record is coherent. It is not
latent in a way that makes it available on demand, because it is not \reachable. It is not
hidden in a way that looking harder would uncover, because exposure is a property of the
interface, not of the observer. It is not merely unobserved, because \occurrent{} records
what happened in the system. And it is not forbidden, because it is \admissible.

The other records are defined in \Cref{supp:certificates}, except the strict-join certificate,
which carries most of the weight of the catalog; we give it in full, filled in for the three-bit
join of \Cref{ex:cube}.

\subsection{The strict-join certificate}
\label{sec:strict-certificate}

\begin{definition}[Strict-join certificate]
\label{def:strict-certificate}
A \emph{strict-join certificate} records a join candidate, composite objecthood,
recoverability of each parent, an anti-product witness, source independence, a settled
budget, a directionality flag and an audit. The
\emph{anti-product witness} has a flag for a joint distinction; flags for factorization on the image
through the left parent, the right parent, the declared product and a declared common
refinement; flags for relabelling-only, scheduling-only and coarsening-only
constructions; a description; and an audit. The witness is valid when the joint-distinction
flag is set, the seven other flags are clear, the description is non-empty and the audit
has an entry. The certificate is valid when the candidate is well formed, the objecthood, retention,
independence and budget flags are set, the witness is valid and the audit has an entry.
Directionality is not required: a strict join does not carry an arrow.
\end{definition}

As a record, the anti-product witness is an assertion of non-factorization: its flags are
entered, not derived. What makes it more than an assertion is \Cref{thm:bridge}: a
semantic strict join that is also proved not to factor, on its image, through a declared
common refinement
and a declared schedule produces a valid witness, with each flag equal to the truth value
of its semantic predicate. Strictness alone does not supply those two proofs
(\Cref{rem:refinement-limit}). The other fields of the certificate (objecthood,
parent retention, source independence and budget) have no semantic model here and remain
recorded assertions.

\begin{example}[The certificate of the three-bit join]
\label{ex:cube-certificate}
For the three-bit join of \Cref{ex:cube} with the observable $f(a,b,h)=h$, take the pairing
$\langle p_A,p_B\rangle$ as the declared common refinement and $p_A$ as the declared schedule.
\Cref{thm:bridge} fills the anti-product witness: the joint-distinction flag is set by the split
pair $(0,0,0)$, $(0,0,1)$, and the seven other flags are clear. A well-formed join candidate,
objecthood of the composite, recoverability of $a$ and $b$, disjoint source roots, a settled
budget and a non-empty certificate audit must then be supplied as recorded assertions; once they are, the certificate is valid by \Cref{def:strict-certificate},
and the directionality flag may stay clear. Only the anti-product witness is derived; the rest is
what a user of the certificate must supply.
\end{example}

